\section{Conclusão}
\label{sec:conclusao}

Este artigo apresentou o AuditEdu-XAI, ferramenta de auditoria algorítmica projetada, desenvolvida e avaliada segundo a DSR, que integra métricas de justiça algorítmica e explicações SHAP em uma arquitetura desacoplada executável em infraestrutura gratuita. A resposta à questão de pesquisa é afirmativa, com as ressalvas apontadas a seguir: a ferramenta detectou, quantificou e explicou vieses de raça/cor e renda em modelos preditivos do ENEM e mostrou-se sensível a vieses aos quais a AUC-ROC pouco respondeu.

A contribuição \textbf{tecnológica} é um artefato, com código publicado em repositório público, que combina métricas de equidade de grupo, decomposição do viés entre dados e modelo, explicações TreeSHAP e uma interface acessível a usuários sem formação em programação; sua implantação pública demonstra que a auditoria algorítmica pode ser oferecida sem custo de infraestrutura, em um padrão potencialmente replicável para outras bases públicas. A contribuição \textbf{empírica} reúne três achados: um modelo de desempenho moderado viola a regra dos quatro quintos para todos os grupos não brancos e para candidatos com renda familiar de até 2,5 salários mínimos, nas três edições; o modelo \textbf{amplifica} a disparidade dos dados, de forma estatisticamente robusta em 13 das 15 combinações de raça/cor e edição e nas faixas de renda mais baixas, levando para baixo do limiar grupos cujo DI nos rótulos estava na fronteira de 0,8; e a ausência de computador na residência é o atributo mais influente, com importância crescente entre 2019 e 2021, em coincidência temporal com o ensino remoto da pandemia, achado compatível com a hipótese de agravamento da exclusão digital, mas que não permite inferência causal. A contribuição \textbf{metodológica} é o protocolo de injeção de viés com controle negativo pareado, que evidenciou que, com 50\% de injeção, a AUC-ROC varia, em média, no máximo 0,032 (Parda, cujo DI cai a quase zero) e menos de 0,001 (Indígena); a avaliação de modelos aplicados a populações heterogêneas não pode, portanto, limitar-se a métricas agregadas.

As principais limitações são: o pequeno número de candidatos indígenas no teste, que torna instáveis suas métricas, em especial a TPR; intervalos de confiança que, com o modelo fixo, não incorporam a variabilidade do treinamento; a auditoria de um único tipo de modelo, que usa raça/cor diretamente e tem hiperparâmetros na borda de uma grade compacta; a natureza associativa dos valores SHAP; o alvo relativo a cada edição, que expressa posição relativa, e não proficiência absoluta; a comparabilidade limitada entre edições; a ausência de avaliação da interface com usuários do público-alvo; a imputação pela moda, que em 2020 atribui um perfil presumido a participantes cuja ausência de resposta provavelmente não é aleatória; e a ausência de medição do tempo de resposta e do consumo de memória em produção.

Como trabalhos futuros, propõem-se: um \textit{bootstrap} com retreinamento do modelo \cite{efron1993bootstrap} e a ampliação da amostra dos grupos minoritários; uma categoria explícita de ``não resposta'' no lugar da moda; a medição do desempenho em produção; a auditoria de um modelo sem raça/cor, para medir a disparidade transmitida por \textit{proxies} \cite{dwork2012fairness}; a integração de técnicas de mitigação, como a reponderação de instâncias e limiares por grupo \cite{hardt2016equality}; a análise de subgrupos interseccionais; a extensão a outros modelos, definições de alvo, edições do ENEM e bases como o SAEB e o Censo Escolar; e testes de usabilidade com o público-alvo.

A regulação brasileira caminha para exigir de sistemas de alto risco avaliações de impacto algorítmico e mecanismos de explicação \cite{brasil2018lgpd, brasil2023pl2338}. Os resultados reforçam essa exigência: um modelo pode parecer adequado por todas as métricas usuais e, ainda assim, excluir sistematicamente os grupos que mais dependem de políticas educacionais justas. A ferramenta torna essa exclusão visível, mensurável e explicável.
